\documentclass[10pt,a4paper]{article}
\usepackage{amsmath,amssymb,amsthm,graphicx}
\usepackage{fontspec}
\setmainfont{DejaVu Serif}
\usepackage{amsmath,amssymb,amsthm,graphicx,xcolor,multirow,mathdots}
\usepackage[shortlabels]{enumitem}
\usepackage[all,2cell]{xy}
\usepackage{tikz,tikz-cd}
\usepackage[unicode,hidelinks]{hyperref}
\usepackage[a4paper,margin=18mm]{geometry}
\numberwithin{equation}{section}
\newenvironment{pf}{\begin{proof}}{\end{proof}}
\newtheorem{theorem}{Theorem}[section]
\newtheorem{lemma}[theorem]{Lemma}
\newtheorem{corollary}[theorem]{Corollary}
\newtheorem{proposition}[theorem]{Proposition}
\newtheorem{conjecture}[theorem]{Conjecture}
\newtheorem{claim}[theorem]{Claim}
\theoremstyle{definition}
\newtheorem{defs}[theorem]{Definitions}
\newtheorem{definition}[theorem]{Definition}
\newtheorem{notation}[theorem]{Notation}
\newtheorem{setup}[theorem]{Setup}
\newtheorem{construction}[theorem]{Construction}
\newtheorem{assumption}[theorem]{Assumption}
\theoremstyle{remark}
\newtheorem{remark}[theorem]{Remark}
\newtheorem{example}[theorem]{Example}
\newtheorem{question}{Question}
\graphicspath{{items/}}
\providecommand{\etalchar}[1]{#1}
\providecommand{\bysame}{\rule[.5ex]{3em}{.3pt}\,}
\providecommand{\doi}[2][]{\href{https://doi.org/#2}{\nolinkurl{doi:#2}}}
\def\Div{\operatorname{Div}}
\def\div{\operatorname{div}}
\def\CDiv{\operatorname{CDiv}}
\def\Cl{\operatorname{Cl}}
\DeclareMathOperator{\Spec}{Spec}
\DeclareMathOperator{\Spf}{Spf}
\DeclareMathOperator{\HH}{H}
\def\Sym{\operatorname{Sym}}
\def\sing{\operatorname{sing}}
\def\mult{\operatorname{mult}}
\def\sym{\operatorname{Sym}}
\def\ex{\operatorname{exp}}

\def\Cox{\operatorname{Cox}}
\def\conv{\operatorname{conv}}
\def\inv{\operatorname{Inv}}
\def\cone{\operatorname{Cone}}
\def\pt{\operatorname{pt}}
\def\lcm{\operatorname{lcm}}

\newcommand{\fromto}[2]{#1, \dotsc, #2}
\newcommand{\setfromto}[2]{\set{\fromto{#1}{#2}}}

\newcommand\TT{{\mathbb{T}}}
\newcommand\CC{{\mathbb{C}}}
\newcommand\PP{{\mathbb{P}}}
\newcommand\RR{{\mathbb{R}}}
\newcommand\QQ{{\mathbb{Q}}}
\newcommand\ZZ{{\mathbb{Z}}}
\newcommand\BB{{\mathbb{B}}}
\newcommand\DDD{{\mathbb{D}}}
\newcommand\SSS{{\mathbb{S}}}


\newcommand\AAA{{\mathbb{A}}}
\newcommand\OO{{\mathbb{O}}}
\newcommand\HHH{{\mathbb{H}}}
\newcommand\EE{{\mathbb{E}}}
\newcommand\FF{{\mathbb{F}}}


\def\A{{\mathbb A}}
\def\B{{\mathbb B}}
\def\C{{\mathbb C}}
\def\F{{\mathbb F}}
\def\G{{\mathbb G}}
\def\H{{\mathbb H}}
\def\N{{\mathbb N}}
\def\Oc{{\mathbb O}}
\def\P{{\mathbb P}}
\def\Q{{\mathbb Q}}
\def\R{{\mathbb R}}
\def\Z{{\mathbb Z}}


\def\cA{{\mathcal A}}
\def\cB{{\mathcal B}}
\def\cC{{\mathcal C}}
\def\cD{{\mathcal D}}
\def\cE{{\mathcal E}}
\def\cF{{\mathcal F}}
\def\cG{{\mathcal{G}}}
\def\cH{{\mathcal{H}}}
\def\cI{{\mathcal I}}
\def\cJ{{\mathcal J}}
\def\cK{{\mathcal K}}
\def\cL{{\mathcal L}}
\def\cM{{\mathcal M}}
\def\cN{{\mathcal{N}}}
\def\cO{{\mathcal{O}}}
\def\cP{{\mathcal{P}}}
\def\cQ{{\mathcal{Q}}}
\def\cR{{\mathcal R}}
\def\cS{{\mathcal S}}
\def\cT{{\mathcal T}}
\def\cU{{\mathcal U}}
\def\cV{{\mathcal V}}
\def\cW{{\mathcal W}}
\def\cX{{\mathcal X}}
\def\cY{{\mathcal Y}}
\def\Q{{\mathbb{Q}}}
\def\G{{\mathbb{G}}}
\def\U{{\mathcal{U}}}
\def\K{{\mathcal{K}}}


\def\fm{{\mathfrak m}}
\def\fn{{\mathfrak n}}
\def\fg{{\mathfrak g}}
\def\fh{{\mathfrak h}}
\def\fp{{\mathfrak p}}
\def\fb{{\mathfrak b}}
\def\fd{{\mathfrak d}}
\def\fe{{\mathfrak e}}
\def\ff{{\mathfrak f}}
\def\fu{{\mathfrak u}}
\def\fsl{{\mathfrak sl}}
\def\fso{{\mathfrak so}}
\def\fsp{{\mathfrak sp}}
\def\sl{\operatorname{SL}}
\def\so{\operatorname{SO}}


\def\cExt{{\mathcal E}xt}
\def\cHom{{\mathcal H}om}

\def\operatorname#1{\mathop{\rm #1}\nolimits}

\def\DA{{\rm A}}
\def\DB{{\rm B}}
\def\DC{{\rm C}}
\def\DD{{\rm D}}
\def\DE{{\rm E}}
\def\DF{{\rm F}}
\def\DG{{\rm G}}


\def\Proj{\operatorname{Proj}}
\def\Aut{\operatorname{Aut}}
\def\CAut{\operatorname{CAut}}
\def\Bl{\operatorname{Bl}}
\def\Chow{\operatorname{Chow}}
\def\Ext{\operatorname{Ext}}
\def\Exc{\operatorname{Exc}}
\def\Hilb{\operatorname{Hilb}}
\def\Hom{\operatorname{Hom}}
\def\Loc{\operatorname{Locus}}
\def\Cloc{\operatorname{Chlocus}}
\def\Mor{\operatorname{Mor}}
\def\Pic{\operatorname{Pic}}
\def\Hom{\operatorname{Hom}}
\def\Sing{\operatorname{Sing}}
\def\reg{\operatorname{Reg}}
\def\Spec{\operatorname{Spec}}
\def\Symm{\operatorname{Sym}}
\def\Tor{\operatorname{Tor}}
\def\codim{\operatorname{codim}}
\def\diag{\operatorname{diag}}
\def\id{\operatorname{id}}
\def\im{\operatorname{im}}
\def\Int{\operatorname{Int}}
\def\length{\operatorname{length}}
\def\rank{\operatorname{rank}}
\def\rk{\operatorname{rk}}
\def\supp{\operatorname{supp}}
\def\red{\operatorname{red}}
\def\deg{\operatorname{deg}}
\def\Bs{\operatorname{Bs}}
\def\Bss{\mathcal{B}}
\def\det{\operatorname{det}}
\def\conj{\operatorname{conj}}
\def\coker{\operatorname{coker}}
\def\gen{\operatorname{gen}}
\def\alg{\operatorname{alg}}
\def\loc{\operatorname{Locus}}
\def\rat{\operatorname{RatCurves}}
\def\pto{{\divideontimes}}
\def\sga{{\langle}}
\def\sgc{{\rangle}}
\def\ojo{{\sphericalangle}}
\def\NE{{\operatorname{NE}}}
\def\ME{{\operatorname{ME}}}
\def\Nef{{\operatorname{Nef}}}
\def\Amp{{\operatorname{Amp}}}
\def\Nu{{\operatorname{N_1}}}
\def\NU{{\operatorname{N^1}}}
\def\Ch{{\operatorname{Ch}}}
\def\Eff{{\operatorname{Eff}}}
\def\Mov{{\operatorname{Mov}}}
\newcommand{\cMov}[1]{\overline{\Mov(#1)}}
\def\Ker{{\operatorname{Ker}}}
\newcommand{\cNE}[1]{\overline{\NE(#1)}}
\newcommand{\nbh}[1]{\operatorname{ngb}(#1)}
\newcommand{\cME}[1]{\overline{\ME(#1)}}
\newcommand{\cEff}[1]{\overline{\Eff(#1)}}
\def\ev{{\operatorname{ev}}}
\def\arg{{\operatorname{arg}}}
\def\rel{{\operatorname{rel}}}
\def\tr{\operatorname{tr}}
\def\Ad{\operatorname{Ad}}
\def\ad{\operatorname{ad}}
\def\Lie{\operatorname{Lie}}
\def\GL{\operatorname{GL}}
\def\PGL{\operatorname{PGL}}
\def\SL{\operatorname{SL}}
\def\Hex{\operatorname{Hex}}
\def\O{\operatorname{\,\!O}}
\def\PSO{\operatorname{PSO}}
\def\SO{\operatorname{SO}}
\def\SP{\operatorname{Sp}}
\def\PSP{\operatorname{PSp}}
\def\Spin{\operatorname{Spin}}
\def\Sec{\operatorname{Sec}}
\def\End{\operatorname{End}}
\def\het{\operatorname{ht}}
%\def\No{\operatorname{N}}
\DeclareMathOperator{\No}{N}
\def\len{\operatorname{\lambda}}
\def\bor{\lhd}

%%%
% PRUNING MAPS:
\def\prul{s}
\def\prur{d}
\newcommand{\orbit}{\curvearrowright}
%%%


%CHOW and all that
\def\Ch{\operatorname{Ch}}
\def\Un{\mathcal{U}}
\def\SX{\mathcal{S}\!X\hspace{-0.5pt}}
\def\GX{\mathcal{G}\!X\hspace{-0.5pt}}
\def\CX{\mathcal{C}\!X\hspace{-0.5pt}}
\def\HX{\mathcal{H}\!X\hspace{-0.5pt}}
\def\ILX{\mathcal{L}\!X\hspace{-0.5pt}}

\def\ss{\operatorname{ss}}
%\renewcommand{\ss}{\operatorname{ss}}



\newcommand{\Chi}{\ensuremath \raisebox{2pt}{$\chi$}}
\newcommand{\pb}{\ar@{}[dr]|{\text{\pigpenfont J}}}
\def\ol{\overline}
\def\tl{\widetilde}
\newcommand{\shse}[3]{0 ~\lra ~#1~ \lra ~#2~ \lra ~#3~ \lra~ 0}
\makeatletter
\newcommand{\xleftrightarrow}[2][]{\ext@arrow 3359\leftrightarrowfill@{#1}{#2}}
\makeatother
\newcommand{\xdasharrow}[2][->]{
\tikz[baseline=-\the\dimexpr\fontdimen22\textfont2\relax]{
\node[anchor=south,font=\scriptsize, inner ysep=1.5pt,outer xsep=2.2pt](x){#2};
\draw[shorten <=3.4pt,shorten >=3.4pt,dashed,#1](x.south west)--(x.south east);
}}
\newcommand{\hooklongrightarrow}{\lhook\joinrel\longrightarrow}
\newcommand{\hooklongleftarrow}{\longleftarrow\joinrel\rhook}

% symmetric difference
%\usepackage{stackengine}
%\def\stacktype{L}
%\stackMath
%\def\dysym{-3pt}
%\newcommand\symdif{\mathrel{\stackon[\dysym]{-}{\dot{\phantom{-}}}}}
\providecommand{\symdif}{\mathbin{\mathpalette\xdotminus\relax}}
\newcommand{\xdotminus}[2]{%
  \ooalign{\hidewidth$\vcenter{\hbox{$#1\dot{}$}}$\hidewidth\cr$#1-$\cr}%
}


\newcommand\m{{\mathfrak m}}
\newcommand\g{{\mathfrak g}}
\newcommand\h{{\mathfrak h}}

\newcommand\ra{{\ \rightarrow\ }}
\newcommand\lra{\longrightarrow}
\newcommand\lla{\longleftarrow}
\newcommand\iso{{\ \cong\ }}

%hats

\newcommand\hX{{\widehat{X}}}
\newcommand\hY{{\widehat{Y}}}
\newcommand\hZ{{\widehat{Z}}}
\newcommand\hL{{\widehat{L}}}
\newcommand\hE{{\widehat{E}}}
\newcommand\hF{{\widehat{F}}}
\newcommand\hPi{{\widehat{\Pi}}}

%% big wedge
\newcommand{\Wedge}[1]{{\textstyle{\bigwedge\nolimits}^{\! #1}}}

%extra symbols

\def\Mo{\operatorname{\hspace{0cm}M}}
\def\Na{\operatorname{\hspace{0cm}M^*}}
\def\We{\operatorname{\hspace{0cm}W}}

\newcommand{\supth}[1]{\ensuremath{#1^{\mathrm{th}}}}
\title{Simple normal-crossing boundary of normalized Chow quotients of convex Picard-one varieties}
\author{Gianluca Occhetta, Luis E. Solá Conde}
\date{arxiv:2503.19663v4; CC BY-SA4.0 current fixed authorv4 and officialPDF; CC BY4.0 separately attributed supporting prior authorv3}
\begin{document}
\maketitle
\setcounter{section}{1}
\setcounter{theorem}{0}
\noindent\textit{Editorial scope.} Count one principal boundary theorem: Theorem 1.1 in its precise form Theorem 3.12. Do not separately count quotient smoothness, Ext propositions, the normal-crossing corollary or other support. The independent cone theorem requires additional primal-curve, intersection and global-contraction obligations.\par\medskip
\noindent\textit{Editorial note.} All counted mathematical text is from fixed current author v4. Its complete current source core is retained, with all supporting statements uncounted. No generated proof or mathematical correction is supplied.\par\medskip
\noindent\textit{Editorial note.} Current author source and original PDFs are CC BY-SA 4.0. Retain attribution and ShareAlike. Prior fixed author v3 is separately CC BY 4.0; the formal NC-ND version is not used as editable source.\par\medskip
\noindent\textit{Editorial note.} Prior labels and citation keys have mechanical namespace prefixes. Statement and proof words are unchanged. Editorial controls preserve original theorem numbering and show the complete prior diagram at a smaller font. Original native files and PDFs are also retained.\par\medskip
\noindent\textit{Editorial note.} The earlier OS26 results are precisely attributed external published inputs, not falsely labeled elementary prerequisites. No unlicensed full source from that work is redistributed.\par\medskip
\noindent\textit{Editorial note.} Count only the principal claim identified here. Final deduplication admission and actual assembly and compilation remain with root.\par\medskip
\noindent\textit{Editorial scope.} Count one principal boundary theorem: Theorem 1.1 in its precise form Theorem 3.12. Do not separately count quotient smoothness, Ext propositions, the normal-crossing corollary or other support. The independent cone theorem requires additional primal-curve, intersection and global-contraction obligations.\par\medskip
\noindent\textit{Allowed standard prerequisites (editorial).}\begin{itemize}
\item Białynicki-Birula local/decomposition theory and equivariant vectorbundle weights. Current§2.1 andBB73,Car02,Ive72; exact normalbundle/cell notation retained.. Established torus-action tools; none of current-specific boundary/primal/cone proofs is removed.
\item Deformation theory of nodal curves, local-to-globalExt, Hilberttangent/obstruction identification. CurrentProposition3.7 proof andTal09§6,Sernesi2006; standard quoted statements retained exactly.. Classical deformation-theoretic inputs; current Kodaira–Spencer surjectivity, torusweight-zero passage and boundary identification are full author text.
\item Numerical divisor/curve duality, nef pullbacks, projectivebundle/blowup Picard behavior, Mori bend-and-break. Current§4; earlier separately licensed amplenessproof refers DebarreProposition3.2.. Established general geometry; current primal construction and its full localconic/globalcontraction argument are not prerequisites.
\end{itemize}\medskip
\section{Background material}

Throughout the paper, we will work over the field of complex numbers. The term {\em rational homogeneous variety} will always refers to a projective variety that supports a transitive algebraic action of a semisimple group. When working with these varieties, we will use the notation of marked Dynkin diagrams, as described, for instance, in \cite[Section~3.1]{WORS5}.

%%%%%%%%%%%%%%%%%%%%%%%%%%%%%%%%%%%%%%%%%%%%%%%%%%%%%%%%%%%
%%%%%%%%%%%%%%%%%%%%%%%%%%%%%%%%%%%%%%%%%%%%%%%%%%%%%%%%%%%
\subsection{$\boldsymbol{\C^*}$-actions}\label{sec:c*}

In this section, we  present some background material on $\C^*$-actions. We refer to \cite{BB,CARRELL,WORS1} for further details. In what follows, $X$ will be a smooth projective variety, endowed with a nontrivial $\C^*$-action.

%%%%%%%%%%%%%%%%%%%%%%%%%%%%%%%%%%%%%%%%%%%%%%%%%%%%%%%%%%%
We denote by $X^{\C^*}$  the fixed locus of the action, and by $\cY$ the set of irreducible fixed-point components: 
$$X^{\C^*}=\bigsqcup_{Y\in \cY}Y.$$ 
 The {\em sink} and the {\em source} of the action, often called the {\em extremal fixed-point components}, are the fixed-point components containing, respectively, the limiting points 
$$
x_-:=\lim_{t\to 0}t^{-1}x, \quad x_+:=\lim_{t\to 0}tx
$$
of a general $x \in X$.  The fixed-point components different from the sink and the source are called {\em inner}.  Every $Y\in\cY$ is smooth (\textit{cf.}~\cite{IVERSEN}), and the normal bundle of $Y$ in $X$ splits into two subbundles, on which $\C^*$ acts with positive and negative weights, respectively:
\begin{equation}
N_{Y|X}\simeq N^+(Y)\oplus N^-(Y).
\label{eq:normal+-}
\end{equation}
We set $\nu^{\pm}(Y):=\rk (N^{\pm}(Y))$.

%%%%%%%%%%%%%%%%%%%%%%%%%%%%%%%%%%%%%%%%%%%%%%%%%%%%%%%%%%%
\subsubsection{Bia{\l}ynicki-Birula decomposition}
Given any set $S \subset X$, we define
\begin{equation}\label{eq:BBcells}
X^\pm(S):=\{x\in X\mid x_\pm\in S\}, \quad
B^\pm(S):=\overline{X^\pm(S)}.
\end{equation}
When $S=Y$ with $Y\in \cY$, we call $X^\pm(Y)$ the positive and the negative {\it Bia{\l}ynicki-Birula cells} of the action at~$Y$  (BB-cells, for short);  we have decompositions 
$$X=\bigsqcup_{Y\in\cY}X^+(Y)=\bigsqcup_{Y\in\cY}X^-(Y).$$  
We will denote by $B^\pm(Y)$ the closures of the BB-cells $X^\pm(Y)$. Every cell $X^\pm(Y)$ has a natural morphism onto~$Y$, sending every $x\in X^\pm(Y)$ to its limiting point $x_\pm$. It was shown by Bia{\l}ynicki-Birula that $X^\pm(Y)$ is an affine bundle, locally isomorphic to $N^\pm(Y)$.

%%%%%%%%%%%%%%%%%%%%%%%%%%%%%%%%%%%%%%%%%%%%%%%%%%%%%%%%%%%
\subsubsection{Actions on polarized pairs}\label{ssec:polar} Given  a line bundle $L\in \Pic(X)$, there exists a linearization of the $\C^*$-action on it, so that for every $Y\in \cY$, $\C^*$ acts on $L_{|Y}$ by multiplication with a character $m\in \Mo(\C^*)=\Hom(\C^*,\C^*)$, called the {\em weight of the linearization on $Y$}. Two linearizations differ by a character of $\C^*$. In particular, for any  $L \in \Pic(X)$, there exists a unique linearization (called {\em normalized linearization})  whose weight at the sink is equal to zero. 
Fix an isomorphism $\Mo(\C^*)\simeq \Z$; then  this linearization defines a map  $\mu_L\colon\cY\to\Z$, called the {\em weight map}. 


A {\em $\C^*$-action on the polarized pair $(X,L)$} is a $\C^*$-action on $X$, together with the weight map $\mu_L$ determined by an ample line bundle $L$.  In this case, we denote by
$$
0=a_0<\dots<a_r
$$ 
the weights $\mu_L(Y)$, $Y\in\cY$, 
and set 
$$
Y_i:=\bigcup_{\mu_L(Y)=a_i}Y.
$$ 
The values $a_i$ will be called the {\em critical values}, and the number $r$ will be called the {\em criticality} of the action.
The minimum and maximum of the critical values are achieved only at the  sink and the  source, 
which are then equal to $Y_0$ and $Y_r$, respectively. The value $\delta=a_r$ is called the {\em bandwidth} of the $\C^*$-action on $(X,L)$. 

%%%%%%%%%%%%%%%%%%%%%%%%%%%%%%%%%%%%%%%%%%%%%%%%%%%%%%%%%%
\subsubsection{Equalized actions}\label{sssec:equal}

A $\C^*$-action  is {\em  equalized} at $Y\in\cY$ if for every $x\in (X^{-}(Y)\cup X^{+}(Y))\setminus Y$, the isotropy group of the action at $x$ is trivial. Equivalently (see \cite[Lemma 2.1]{WORS3}), $\C^*$ acts on $N^+(Y)$ with all the weights equal to $+1$ and on $N^-(Y)$ with all the weights equal to $-1$. The action is called {\em equalized} if it is equalized at every $Y \in \cY$.

For an equalized action, the closure of any $1$-dim\-ens\-ional orbit is a smooth rational curve, whose $L$-degree may be computed in terms of the weights at its extremal points (\textit{cf.} \cite[Lemma 2.2]{WORS3}, \cite[Corollary 3.2]{RW}).  

\begin{lemma}[AM vs. FM]\label{lem:AMvsFM}
Let $(X,L)$ be a polarized pair with an equalized $\C^*$-action, and let $C$ be the closure of a $1$-dimensional orbit, with sink  $x_-$ and source $x_+$. Then $C$ is a smooth rational curve of\, $L$-degree  equal to $\mu_L(x_+)-\mu_L(x_-)$.
\end{lemma} 

Using Lemma~\ref{lem:AMvsFM}, one can show (see \cite[Corollary 2.4]{OS2}) the following. 

\begin{corollary} \label{cor:can} 
Let $X$ be a smooth projective variety  with an equalized $\C^*$-action, and let $C$ be the closure of a $1$-dimensional orbit, with sink  in $Y$ and source in  $Y'$. Then 
\[
-K_X \cdot C =\left(\nu^+(Y')-\nu^-(Y')\right)-\left(\nu^+(Y) -\nu^-(Y)\right).
\]
\end{corollary}

%%%%%%%%%%%%%%%%%%%%%%%%%%%%%%%%%%%%%%%%%%%%%%%%%%%%%%%%%%%
%%%%%%%%%%%%%%%%%%%%%%%%%%%%%%%%%%%%%%%%%%%%%%%%%%%%%%%%%%%
\subsection{Quotients}\label{sssec:quotients} For  more details on the contents of this section,  we refer the reader to \cite{WORS3} and \cite{WORS6}. A $\C^*$-action on a polarized variety $(X,L)$ allows one to define projective quotients of $X$, as follows. Set, for every $\tau \in [0,\delta]\cap \Q$, 
 \[A(\tau):=\bigoplus_{\substack{m\geq 0\\m\tau\in\Z}}\HH^0(X,mL)_{m\tau},\quad \GX(\tau):=\Proj(A(\tau)).
\]
By \cite[Section 2.2]{WORS6}, $\GX(\tau)$ is the same for any  $\tau\in (a_i,a_{i+1})\cap \Q$, so we set 
\begin{itemize}[itemsep=5pt,topsep=5pt]
\item $\GX_{i,i}:=\GX(\tau)$ for $\tau=a_i$,  
\item $\GX_{i,i+1}:=\GX(\tau)$ for $\tau\in(a_i,a_{i+1})\cap \Q$.
\end{itemize}
Note that $\GX_{i,i+1}$ is a {\em geometric quotient} of $X$ for every $i$, while $\GX_{i,i}$ is only a good categorical quotient (see \cite[Section 6.1, Definition, p.~94]{DolgachevGIT}). Following \cite[Chapter 3, p.~10]{BBC}, we say that $\GX_{i,i}$ is {\em semigeometric}.  
Furthermore, the varieties $\GX_{i,i+1}$ are birationally equivalent, and the natural maps among them fit in the following commutative diagram, where the diagonal arrows are contractions and the horizontal maps are Atiyah flips (in the sense of \cite[Definition 6.1]{WORS1}): 
\begin{equation}
\begin{tikzcd}[
  column sep={3.4em,between origins},
  row sep={2.7em,between origins},
]
&\GX_{0,1} \arrow[rr,dashed] \arrow[rd] \arrow[dl]&&\GX_{1,2} \arrow[dl]&\dots\ \ &\GX_{r-2,r-1} \arrow[rr,dashed] \arrow[rd] &&\GX_{r-1,r}\arrow[rd] \arrow[dl]&\\
\GX_{0,0}&&\GX_{1,1}&\dots&&\dots\ \ &\GX_{r-1,r-1}&&\GX_{r,r}\rlap{.}
\end{tikzcd}
\label{eq:GITquot}
\end{equation}

%%%%%%%%%%%%%%%%%%%%%%%%%%%%%%%%%%%%%%%%%%%%%%%%%%%%%%%%%%
\subsubsection{Prunings} \label{ssec:prun}

Given an equalized $\C^*$-action of criticality $r$ and bandwidth $\delta$ on the smooth polarized pair $(X,L)$,  we consider the blowup $\beta\colon X^\flat\to X$ of $X$ along the sink $Y_0$  and the source $Y_r$, with exceptional divisors $Y^\flat_0,Y^\flat_r$, and the induced $\C^*$-action.
Given rational numbers $\tau_-<\tau_+$ in $[0,\delta]$, we set
\[
L(\tau_-,\tau_+):=\beta^*L-\tau_-Y^\flat_0-(\delta-\tau_+)Y^\flat_r 
\]
and define the {\em pruning} of $(X,L)$ with respect to $\tau_-,\tau_+$ as
\[X(\tau_-,\tau_+):=\Proj \bigoplus_{m\in k\Z_{\geq 0}}\HH^0(X^\flat,mL(\tau_-,\tau_+)),\]
where $k>0$ is the minimum integer such that $k\tau_\pm\in \Z_{> 0}$.
As shown in \cite[Lemma 3.5]{WORS6}, for values $\tau_-,\tau_+\in([0,\delta]\cap\Q)\setminus\{a_0,\dots a_r\}$, the pruning $X(\tau_-,\tau_+)$ depends only on the intervals $(a_i,a_{i+1})$, $ (a_{j-1},a_j)$ containing $\tau_-,\tau_+$, respectively. Hence we can set
$$X_{i,j}:=X(\tau_-,\tau_+)\quad\mbox{whenever } \tau_-\in(a_i,a_{i+1}), \;\tau_+\in (a_{j-1},a_j).$$ 

%%%%%%%%%%%%%%%%%%%%%%%%%%%%%%%%%%%%%%%%%%%%%%%%%%%%%%%%%%
\subsubsection{Chow quotient}

Given a projective variety $X$ with a nontrivial $\C^*$-action, the {\em normalized Chow quotient} of the action is the normalization $\CX$ of the closure of the subscheme of $\Chow(X)$ parametrizing the closures of general orbits.

More precisely, one may find a nonempty open set $V\subset X$ of points $x$ whose orbits $\overline{\C^*\cdot x}$ belong to a single homology class in $X$. By shrinking $V$, if necessary, the geometric quotient of $V$ by the action of $\C^*$ exists. Taking closures in $X$ of the orbits of points of $V$, we obtain an algebraic irreducible family of rational curves in $X$ and a morphism $\psi\colon V/\C^*\to \Chow(X)$.

The {\em Chow quotient} of $X$ by the action of $\C^*$ is then defined as the closure 
$$
\overline{\CX}:=\overline{\psi(V/\C^*)}\subset \Chow(X).
$$
Its normalization $\CX$ will be called the {\em normalized Chow quotient} of $X$ by the $\C^*$-action.  Pulling back the universal family of $\Chow(X)$ to $\CX$, we obtain a family of $1$-dimensional varieties $p\colon\cU\to \CX$, with evaluation morphism $q\colon\cU\to X$.

\begin{remark}\label{rem:chains} 
Let us consider a smooth polarized pair $(X,L)$ endowed with an equalized $\C^*$-action of criticality~$r$, with critical values $a_0,\dots,a_r$. By \cite[Lemma~4.3]{WORS6}, the fibers of $p$ are reduced $1$-dimensional schemes with smooth irreducible components, whose dual graphs are of type $\DA$. Thus any element of the family parametrized by $\CX$ takes the form $C=C_1\cup\dots\cup C_k$, 
where 
\begin{enumerate}[label=(C\arabic*),ref=C\arabic*]
\item\label{C1} every $C_j$ is the (smooth) closure of a $1$-dimensional orbit; 
\item\label{C2} the singular points of $C$ are the (transversal) intersections $P_j:=C_j\cap C_{j+1}$ and are fixed points of weights $a_{i_j}$, with  $0<i_1<\dots<i_{k-1}<r$; 
\item\label{C3} every $P_j$ is the source of $C_j$ and the sink of $C_{j+1}$, for $j=1,\dots,k-1$;  
\item\label{C4} the sink of $C_1$ and the source of $C_k$ belong to $Y_0$ and $Y_r$, respectively.
\end{enumerate}
Note that property~(\ref{C4}) follows from Lemma~\ref{lem:AMvsFM} together with the fact that the $L$-degree of $C$ equals the $L$-degree of the general element of $\CX$, which is $\delta=a_r-a_0$.
\end{remark}

\begin{remark}\label{rem:Hilb} From this description of the fibers of $p$, one may prove (see \cite[Lemma~4.5]{WORS6}) that the family $p$ is flat. In particular, $\CX$ is isomorphic to the normalization of an irreducible $\C^*$-invariant subvariety of $\Hilb(X)$, which we call the {\em Hilbert quotient} of the action (see \cite[Remark~4.6]{WORS6}). 
\end{remark}

\textit{A priori}, it is not clear that any $\C^*$-invariant subscheme $C=C_1\cup\dots\cup C_k$ satisfying the properties (\ref{C1})--(\ref{C4}) listed in Remark~\ref{rem:chains} belongs to $\CX$. However, this has been proved to be true, under the following mild additional assumptions. 

\begin{setup}\label{set:WORS6}
The pair $(X,L)$ is a smooth polarized pair endowed with an equalized $\C^*$-action of criticality $r$, with critical values $0=a_0<\dots<a_r$, such that $\nu^\pm(Y) \ge 2$ for every fixed-point component of weight $a_i$, with $2 \le i \le r-2$.
\end{setup} 

\begin{remark}\label{rem:picone}
The technical assumption above on the $\nu^\pm(Y)$ is fulfilled if the variety $X$ has Picard number~$1$ (\textit{cf.}~\cite[Lemma 2.8]{WORS1} and \cite[Lemma 4.5]{WORS3}).
\end{remark}


\begin{corollary}[\textit{cf.} {\cite[Corollary~5.9]{WORS6}}]\label{cor:CXcontainsall}
Let $(X,L)$ be as in Setup~\ref{set:WORS6}. Then every $\C^*$-invariant $1$\nobreakdash-di\-men\-sional reduced subscheme $C$ of\, $X$ with irreducible components $C_1,\dots,C_k$ satisfying the properties {\rm(\ref{C1})--(\ref{C4})}  listed in Remark~\ref{rem:chains}, is an element of the family parametrized by $\CX$.
\end{corollary}

The proof  uses the recursive description of $\CX$ given in \cite[Theorem 1.1 and Corollary 7.1]{WORS6}, which we now recall. 

\begin{theorem}\label{thm:rhombuses}
Let $(X,L)$ be as in Setup~\ref{set:WORS6}. Denote by $X_{i,j}$, $0\leq i<j\leq r$, the corresponding prunings of\, $X$, and by $\CX_{i,j}$ their normalized Chow quotients. Then $\CX=\CX_{0,r}$, and we have a commutative diagram 
\begin{equation}\label{fig:Chow0}
\begin{tikzcd}[
  column sep={2.45em,between origins},
  row sep={2.7em,between origins}]
&&&&&&\CX_{0,r}\arrow[rd,"\prur"] \arrow[dl,"\prul",labels=above left] &&&&&&\\
&&&&&\CX_{0,r-1}\arrow[rd,"\prur"] \arrow[dl,"\prul",labels=above left]&&\CX_{1,r}\arrow[rd,"\prur"] \arrow[dl,"\prul",labels=above left]&&&&&\\
&&&&\CX_{0,r-2} \arrow[dl,"\prul",labels=above left]\arrow[rd,"\prur"]&&\CX_{1,r-1}\arrow[rd,"\prur"] \arrow[dl,"\prul",labels=above left]&&\CX_{2,r}\arrow[dl,"\prul",labels=above left]\arrow[rd,"\prur"]&&&&\\
&&& \arrow[ld,dotted,no head]&& \arrow[ld,dotted,no head]&&\arrow[rd,dotted,no head]&&\arrow[rd,dotted,no head]&&&\\
&&\CX_{0,2}\arrow[rd,"\prur"] \arrow[dl,"\prul",labels=above left]&&\CX_{1,3} \arrow[dl,"\prul",labels=above left]&&\dots&&\CX_{r-3,r-1}\arrow[rd,"\prur"] &&\CX_{r-2,r}\arrow[rd,"\prur"] \arrow[dl,"\prul",labels=above left]&&\\
&\GX_{0,1}\arrow[ld]\arrow[rd]&&\GX_{1,2}\arrow[ld]\arrow[rd]&&\dots&&\dots&&\GX_{r-2,r-1}\arrow[ld]\arrow[rd]&&\GX_{r-1,r}\arrow[ld]\arrow[rd]&\\
\GX_{0,0}&&\GX_{1,1}&&\GX_{2,2}&&\dots&&\GX_{r-2,r-2}&&\GX_{r-1,r-1}&&\GX_{r,r}\rlap{,}
\end{tikzcd}
\end{equation}
where all the arrows denoted by $\prul,\prur$ are blowups and every rhombus in the diagram is a normalized Cartesian square.
\end{theorem}

In categorical language, we may simply say that $\CX$ is the normalized inverse limit of the GIT quotients
\[\begin{tikzcd}[
  column sep={2.5em,between origins},
  row sep={2.9em,between origins}]
&\GX_{0,1}\arrow[ld]\arrow[rd]&&\GX_{1,2}\arrow[ld]\arrow[rd]&&\dots&&\dots&&\GX_{r-2,r-1}\arrow[ld]\arrow[rd]&&\GX_{r-1,r}\arrow[ld]\arrow[rd]&\\
\GX_{0,0}&&\GX_{1,1}&&\GX_{2,2}&&\dots&&\GX_{r-2,r-2}&&\GX_{r-1,r-1}&&\GX_{r,r}\rlap{.}
\end{tikzcd}\]

The centers of the blowups $\prul,\prur$ appearing in Diagram (\ref{fig:Chow0}) have been described in \cite[Remark~5.10]{WORS6}; they are smooth for the maps appearing in the first layer right above the geometric quotients (see \cite[Lemma~5.7]{WORS6}):
\[\begin{tikzcd}[
  column sep={2.6em,between origins},
  row sep={2.9em,between origins}]
&\CX_{0,2}\arrow[ld,"\prul",labels=above left]\arrow[rd,"\prur"]&
&\CX_{1,3} \arrow[dl,"\prul",labels=above left]&&\dots&&\CX_{r-3,r-1}\arrow[rd,"\prur"] &
&\GX_{r-2,r}\arrow[ld,"\prul",labels=above left]\arrow[rd,"\prur"]&\\
\GX_{0,1}&&\GX_{1,2}&
&\dots&&\dots&
&\GX_{r-2,r-1}&&\GX_{r-1,r}\rlap{.}
\end{tikzcd}\]
For the maps in successive layers, the centers are not necessarily smooth, as shown in \cite[Example~6.1]{WORS6}). However, the situation improves under certain hypotheses. Let us recall the following. 

\begin{definition}\label{def:convex}
A smooth projective variety $X$ is said to be {\em convex} if for every morphism $\mu\colon\P^1\to X$ we have that $\HH^1(\P^1,\mu^*T_X)=0$. 
\end{definition}

The following  has been shown in \cite[Theorem~1.2]{WORS6}. 

\begin{theorem}\label{thm:WORS6smooth}
Let $(X,L)$ be as in Setup~\ref{set:WORS6}. Assume moreover that $X$ is convex. Then the varieties $\CX_{i,j}$ for $i=0$ or $j=r$ are smooth, and the maps 
\[\xymatrix{\GX_{0,1}&\CX_{0,2}\ar[l]_(0.44)s&\cdots\ar[l]_(0.42)s&\CX\ar[l]_(0.42)s\ar[r]^(0.44)d&\cdots\ar[r]^(0.40)d&\CX_{r-2,r}
\ar[r]^(0.46)d&\GX_{r-1,r}}\]
are blowups with smooth centers.
\end{theorem}

\section{Chow quotients of convex varieties of Picard number 1}\label{sec:boundaryprimal}

Throughout the rest of the paper, we will work in the following setting. 

\begin{setup}\label{set:main}
Let $X$ be a smooth convex variety of Picard number $1$, endowed with an equalized nontrivial $\C^*$-action of criticality $r$. 
We will denote by  $\cM$ a family of rational curves of minimal degree on $X$ (which exists since $X$ is uniruled), by $p\colon\cU \to \cM$ its universal family, and by $q\colon \cU \to X$ the evaluation morphism (see \cite[Proposition~II.2.11]{kollar}).
\end{setup}

\begin{remark}\label{rem:convexcurves}
Since $X$ is convex, $\cM$ is {\em beautiful}; \textit{i.e.}, the evaluation morphism $q\colon \cU \to X$ is smooth and surjective (see \cite[Section 2.7]{OS2} for further details). We will denote by $\cU_x$ the fiber of $q$ over the point $x \in X$. Note that since $X$ is a Fano manifold (in particular, simply connected), then  every $\cU_x$ is connected, therefore irreducible (\textit{cf.}~\cite[Proposition~1.6]{Kan5}). 
\end{remark}

Let us recall some results recently obtained in \cite{OS2}. 
The $\C^*$-action on $X$ induces actions on $\cM$ and~$\cU$ such that $p$ and $q$ are equivariant morphisms. These actions are equalized (see \cite[Propositions 4.5 and 4.8]{OS2}); in particular, the closure of every  $1$-dimensional orbit in $\cM$ and $\cU$ is a smooth rational curve.

\begin{remark}\label{rem:f2}
Given a rational curve $\Gamma \subset \cM$ that is the closure of a $1$-dimensional orbit, $S_\Gamma:=p^{-1}(\Gamma)$ is a rational ruled surface $\FF_e$ with $e \le 2$ (see \cite[Proposition 4.6]{OS2}). By the same proposition, if no point of~$\Gamma$ parametrizes a curve that is fixed by the action on $X$, then $e=0$ or $e=2$ and the action has four fixed points (see \cite[Figure 4.1]{OS2}).
\end{remark}
 
 The following statement has been proved in \cite[Theorem 5.3]{OS2}.
 
 \begin{proposition}\label{prop:onecomp}
 In the situation of Setup~\ref{set:main}, denoting by $L$ the ample generator of\, $\Pic(X)$, the $L$-weights of the action are $0 < d < 2d < \dots < (r-1)d < rd = \delta$, and, for every $j =0, \dots, r$, there exists a unique $Y_j \in \cY$ such that $\mu_L(Y_j)=jd$.
\end{proposition}

\begin{remark}\label{rem:convex} 
 Proposition~\ref{prop:onecomp} allows us to use the following notation:
  \[
  B_j^\pm:=B^\pm(Y_j),\;\; \nu^\pm_j:=\nu^\pm(Y_j)\quad\mbox{for every $j=0, \dots, r$}\quad \mbox{and}\quad \nu^+_{-1}=\nu^-_{r+1}:=0.\] Note that, by definition, $\nu^+_0=\nu^-_r=0$, $\nu^-_0=\dim X-\dim Y_0$, and $\nu^+_r=\dim X-\dim Y_r$.
\end{remark}

\subsection{Boundary divisors}\label{ssec:boundary}


The main result of this section is that the subscheme of $\CX$ parametrizing reducible cycles is a union of smooth divisors, which we call {\em boundary divisors}, with simple normal crossings. 

Our first observation is the following. 

\begin{remark}\label{rem:Hilb2}
Since $X$ is a convex variety, $\CX$  is smooth by Theorem~\ref{thm:WORS6smooth}, which has been proven in \cite[Proposition~6.4]{WORS6}. The arguments there, applied directly to $\CX$, show when $\CX$ is considered as a subscheme of $\Hilb(X)$ (see Remark~\ref{rem:Hilb}), the convexity assumption implies that $\Hilb(X)$ is smooth at every point of $\CX$. Then $\CX$, which is the only $\C^*$-fixed-point component of $\Hilb(X)$ containing the class of the closure of the general orbit in $X$, is smooth as well.
\end{remark}

Given a closed point $[C]$ of $\CX\subset \Hilb(X)$, it determines a $\C^*$-invariant subscheme $C$ of $X$. Moreover, as observed in \cite[Proof of Proposition 6.4]{WORS6}, the singularities of $C$ are nodes. Let us now recall some facts about deformations of projective nodal curves. We refer the interested reader to \cite[Section~6]{Talpo} for a detailed account on this topic.

\begin{proposition}\label{prop:deforms}
Let $C$ be a projective connected curve whose irreducible components are smooth rational curves $C_1,\dots, C_k$ and whose singular locus is the set of points $\{P_i=C_{i}\cap C_{i+1}|\,\,i=1,\dots,k-1\}$. We assume that the singularities of\, $C$ at these points are ordinary nodes. Then 
\begin{itemize}[leftmargin=25pt]
\item $\cExt^1(\Omega_C,\cO_C)$ is supported on the singular points $P_i$, $i=1,\dots,k-1$, and $$\cExt^1(\Omega_C,\cO_C)_{P_i}\simeq \C$$
for every $i=1,\dots,k-1$;
\item the complex vector space $\Ext^1(\Omega_C,\cO_C)$ of first-order deformations of $C$  has dimension $k-1$, and it is naturally isomorphic to
\[
\HH^0(C,\cExt^1(\Omega_C,\cO_C))\simeq \bigoplus_{i=1}^{k-1}\cExt^1(\Omega_C,\cO_C)_{P_i}.
\]
\end{itemize}
\end{proposition}

\begin{proof}
For the first assertion, we refer the reader to \cite[Remark~6.4.2]{Talpo}. The second claim follows from the local-to-global
$\Ext$ spectral sequence.
\end{proof}

\begin{remark}\label{rem:deforms}
  Proposition~\ref{prop:deforms} can be interpreted geometrically as follows. The first-order deformations of $C$ correspond to smoothings of $C$ at its singular points (see \cite[Sections~6.3 and~6.4]{Talpo} for details).  More precisely, for every $i$, the $1$-dimensional vector space $\cExt^1(\Omega_C,\cO_C)_{P_i}$ corresponds, in the space $\Ext^1(\Omega_C,\cO_C)$ of first-order deformations of $C$, to deformations that smoothen the point $P_i$ and are locally trivial around the $P_j$ with $j\neq i$.

In this way, when we choose isomorphisms $\cExt^1(\Omega_C,\cO_C)_{P_i}\simeq \C$ and denote by $t_i$, $i=1,\dots , k-1$, the compositions
$t_i\colon\Ext^1(\Omega_C,\cO_C)\to \cExt^1(\Omega_C,\cO_C)_{P_i}\simeq \C$,  the deformations of $C$ that are locally trivial around $P_i$ are given by the equation $t_i=0$.
\end{remark}

We will now go back to the case in which $C$ is the subscheme of $X$ determined by a closed point $[C]\in \CX$ and see that every first-order deformation of  $C$ extends to a deformation of $C$ as a $\C^*$-invariant subscheme of $X$. More precisely, we will show the following. 

\begin{proposition}\label{prop:kodspen}
Let $[C]$ be a closed point of\, $\CX$. Then  the restriction to $T_{\CX,[C]}$ of the Kodaira--Spencer map 
\[
T_{\Hilb(X),[C]}\lra \Ext^1(\Omega_C,\cO_C)
\]
is surjective.
\end{proposition}

\begin{proof}
As observed in \cite[Proof of Proposition 6.4]{WORS6}, the singularities of $C$ are local complete intersections; then we have a short exact sequence
\begin{equation}\label{eq:ses1}
\shse{\cI/\cI^2}{{\Omega_X}_{|C}}{\Omega_C}
\end{equation}
and  $\Ext^1({\Omega_X}_{|C},\cO_C)=0$. 
We get the exact sequence 
\begin{multline}\label{eq:FP}
0\lra\Hom(\Omega_C,\cO_C)\lra\Hom({\Omega_X}_{|C},\cO_C)\lra \Hom(\cI/\cI^2,\cO_C)\lra \Ext^1(\Omega_C,\cO_C)\lra 0.
\end{multline}
The space $\Hom(\cI/\cI^2,\cO_C)$ is the Zariski tangent space of $\Hilb(X)$ at $[C]$, and the map $\Hom(\cI/\cI^2,\cO_C)\to \Ext^1(\Omega_C,\cO_C)$ is the Kodaira--Spencer map, associating  with every first-order deformation of $C$ as a subscheme of $X$ the corresponding first-order deformation of $C$ as a scheme over $\Spec(\C)$ (\textit{cf.}~\cite[Remark 2.4.4]{Ser1}). We then see that this map is surjective. 

Now we add the action of $\C^*$ into the picture. On one hand, we have induced $\C^*$-actions on the  spaces $\Ext^1(\Omega_C,\cO_C)$ and $\HH^0(C,\cExt^1(\Omega_C,\cO_C))$, and the natural map between them is $\C^*$-equivariant. Hence, since $\cExt^1(\Omega_C,\cO_C)$ is supported on the singular points of $C$, we conclude that the action on $\Ext^1(\Omega_C,\cO_C)$ is trivial.

 On the other hand, the tangent space $T_{\CX,[C]}$ is the weight zero eigenspace  $\Hom(\cI/\cI^2,\cO_C)_0\subset \Hom(\cI/\cI^2,\cO_C)$, which must surject onto the weight zero eigenspace $\Ext^1(\Omega_C,\cO_C)_0=\Ext^1(\Omega_C,\cO_C)$. This completes the proof.
\end{proof}

\begin{corollary}\label{cor:kodspen}
In the situation of Setup~\ref{set:main}, the subscheme $\CX_{\sing}$ of\, $\CX$ parametrizing singular $\C^*$-invariant subschemes is a divisor with normal crossings.
\end{corollary}

\begin{proof}
The fact that $\CX_{\sing}$ is a divisor follows from \cite[Propositions 6.2.11 and 6.3.5]{Talpo}.  We will now study its local equations.

Let $[C]$ be a closed point, corresponding to a singular $\C^*$-invariant subscheme $C\subset X$, with irreducible components $C_1,\dots, C_k$ and singular points $P_i=C_i\cap C_{i+1}$, $i=1,\dots,k-1$. Let us  consider again the Zariski tangent space $T_{\CX,[C]}$ and the linear map
$$
\phi\colon T_{\CX,[C]}\lra \Ext^1(\Omega_C,\cO_C)\simeq \bigoplus_i\cExt^1(\Omega_C,\cO_C)_{P_i}
$$ 
defined as the restriction to $T_{\CX,[C]}$ of the Kodaira--Spencer map, which is surjective by Proposition~\ref{prop:kodspen}. We pull back via $\phi$ the linear coordinates $(t_1,\dots, t_{k-1})$ in  $\Ext^1(\Omega_C,\cO_C)$ introduced in Remark~\ref{rem:deforms}; abusing notation, we denote these linear forms by $t_i\colon T_{\CX,[C]}\to\C$.  This tells us that locally around $[C]$ the divisor $\CX_{\sing}\subset\CX$ is given by an equation of the form $\ol{t}_1\cdot\ldots \cdot \ol{t}_{k-1}=0$, where $(d\ol{t}_i)_{[C]}=t_i$.  In particular, it has normal crossings at $[C]$. 
\end{proof}

To finish our description of $\CX_{\sing}$,  we introduce the following divisors.  

\begin{definition}\label{def:Ei}
For  every $i\in \{1,\dots,r-1\}$, we denote by $E^\circ_i\subset\CX$ the set parametrizing cycles having a unique singular point in the component $Y_i$. By Corollary~\ref{cor:CXcontainsall} and the Bia{\l}ynicki-Birula theorem, $E^\circ_i$ is isomorphic to an open subset of $\P(N_{Y_i}^-)\times_{Y_i}\P(N_{Y_i}^+)$, which has dimension $n-1$. We denote by $E_i$ its closure in $\CX$ and call it the $\supth{i}$ {\em boundary divisor of }\, $\CX$. 
\end{definition}



\begin{theorem}\label{thm:smoothborder} 
Let $X$ be a convex variety of Picard number $1$, endowed with an equalized $\C^*$-action of criticality $r\geq 2$. Then the Chow quotient $\CX$ of\, $X$ by the action of\, $\C^*$ is smooth, and the subscheme $\CX_{\sing}\subset\CX$ parametrizing reducible elements is a divisor with simple normal crossings, equal to
$$\sum_{i=1}^{r-1}E_i.$$
Moreover, we have $[C] \in E_i$ if and only if the cycle $C$ has a singular point in $Y_i$. 
\end{theorem}

\begin{remark}\label{rem:main1}
Note that in the case in which the criticality of the action is $r=1$, Theorem~\ref{thm:smoothborder} is obvious. In fact, in this case the Chow quotient $\CX$ equals the geometric quotient $\GX_{0,1}$, which is the projectivization of the normal bundle of $Y_0$ in $X$, hence smooth, and the subscheme of $\CX$ parametrizing reducible $1$-cycles is empty.
\end{remark}

\begin{proof}[Proof of Theorem~\ref{thm:smoothborder}]
Take an element $[C]\in \CX$ consisting of $k$ irreducible components $C_1,\dots,C_k$, and denote by $P_1\in Y_{i_1},\dots, P_{k-1}\in Y_{i_{k-1}}$ its singular points. 

With the notation of the proof of Corollary~\ref{cor:kodspen}, for every $j\in\{1,\dots,k-1\}$, there exists a smooth  curve $\gamma_j$ in $\CX$ passing through
 $[C]$, whose tangent space at $[C]$ is generated by $v\in T_{\CX,[C]}$, such that $\phi(v)\in \Ext^1(\Omega_C,\cO_C)$ satisfies $t_j(\phi(v))=1$ and $t_s(\phi(v))=0$ for $s\neq j$. 

The general element of $\gamma_j$ has a unique singular point, which belongs to $Y_{i_j}$; in fact, the singular points of an element of $\CX$ are fixed points, and the unique singular point of an element of $\gamma_j$ different from $C$ is a deformation of $P_j\in Y_{i_j}$. In particular, the general element of $\gamma_j$ belongs to  $E_{i_j}$, so $[C]\in E_{i_j}$ as well. We conclude that $[C]\in \bigcap_{j=1}^{k-1}E_{i_j}$.


On the other hand, the arguments in the proof of Corollary~\ref{cor:kodspen} tell us that, locally around $[C]$, $\CX_{\sing}$ is the normal crossing intersection of exactly $k-1$ divisors; hence these divisors are necessarily the $E_{i_j}$.  This finishes the proof.
\end{proof}

\clearpage
\section*{Prior author-source supplement (uncounted)}
The following separately attributed excerpts are from Occhetta, Romano, Sol'a Conde and Wi'sniewski, arXiv:2310.18623v3 (CC BY 4.0). They supply the earlier Hilbert-scheme argument explicitly invoked by the current article. No result from this prior work is counted here.
\renewcommand{\theHsection}{wors6.\arabic{section}}
\renewcommand{\theHtheorem}{wors6.\arabic{section}.\arabic{theorem}}
\renewcommand{\theHequation}{wors6.\arabic{section}.\arabic{equation}}\renewcommand{\theequation}{\arabic{equation}}
\setcounter{section}{1}\setcounter{theorem}{0}\setcounter{equation}{1}
\begingroup\small

We explore here this idea, from the viewpoint of birational geometry. All the geometric quotients of the action are birationally dominated by $\CX$; in particular, we get a resolution  of the birational map $\psi:\GX_{0,1}\dashrightarrow \GX_{r-1,r}$:
$$
\xymatrix{&\CX\ar[dr]\ar[dl]&\\\GX_{0,1}\ar@{-->}[rr]^{\psi}&&\GX_{r-1,r}}
$$
It then makes sense to study the properties of this resolution and ask ourselves, for instance, whether it gives a Strong Factorization of $\psi$, that is whether the morphisms $\CX\to \GX_{0,1}$, $\CX\to \GX_{r-1,r}$ are compositions of blowups along smooth centers. What we prove is a structure theorem that shows how to obtain $\CX$ recursively out of the Chow quotients of the birational modifications $X_{i,j}$ described above, via subsequent blowups. More precisely, our main statement is the following (see $\S$\ref{wors6:sssec:quotients} for the definition of $B^\pm(Y)$):

\begin{theorem}\label{wors6:thm:main}
Let $(X,L)$ be a polarized pair, with $X$ smooth, endowed with an equalized nontrivial $\C^*$-action such that  $\codim(B^\pm(Y),X)>1$ for every inner fixed point component $Y$. Denote by $r$ be the criticality of the action, by $X_{i,j}$, $0\leq i<j\leq r$, the corresponding prunings of $X$, and by $\CX_{i,j}$ their normalized Chow quotients. Then $\CX=\CX_{0,r}$, and we have a commutative diagram:
\begin{equation}\label{wors6:fig:Chow0}
\begin{tikzcd}[
  column sep={2.7em,between origins},
  row sep={3.2em,between origins}]
&&&&&&\CX_{0,r}\arrow[rd,"\prur"] \arrow[dl,"\prul",labels=above left] &&&&&&\\
&&&&&\CX_{0,r-1}\arrow[rd,"\prur"] \arrow[dl,"\prul",labels=above left]&&\CX_{1,r}\arrow[rd,"\prur"] \arrow[dl,"\prul",labels=above left]&&&&&\\
&&&&\CX_{0,r-2} \arrow[dl,"\prul",labels=above left]\arrow[rd,"\prur"]&&\CX_{1,r-1}\arrow[rd,"\prur"] \arrow[dl,"\prul",labels=above left]&&\CX_{2,r}\arrow[dl,"\prul",labels=above left]\arrow[rd,"\prur"]&&&&\\
&&& \arrow[ld,dotted,no head]&& \arrow[ld,dotted,no head]&&\arrow[rd,dotted,no head]&&\arrow[rd,dotted,no head]&&&\\
&&\CX_{0,2}\arrow[rd,"\prur"] \arrow[dl,"\prul",labels=above left]&&\CX_{1,3} \arrow[dl,"\prul",labels=above left]&&\dots&&\CX_{r-3,r-1}\arrow[rd,"\prur"] &&\CX_{r-2,r}\arrow[rd,"\prur"] \arrow[dl,"\prul",labels=above left]&&\\
&\GX_{0,1}\arrow[rd]&&\GX_{1,2}\arrow[ld]\arrow[rd]&&\dots&&\dots&&\GX_{r-2,r-1}\arrow[ld]\arrow[rd]&&\GX_{r-1,r}\arrow[ld]&\\
&&\GX_{1,1}&&\GX_{2,2}&&\dots&&\GX_{r-2,r-2}&&\GX_{r-1,r-1}&&
\end{tikzcd}
\end{equation}
where all the arrows denoted by $\prul,\prur$ are blowups and every rhombus in the diagram is a normalized Cartesian square.
\end{theorem}

\endgroup

\begin{theorem}\label{wors6:thm:convex}
Under the assumptions of Theorem \ref{wors6:thm:main}, assume moreover that $X$ is convex. Then the varieties $\CX_{i,j}$ with $i=0$ or $j=r$ are smooth, and the maps 
\[\xymatrix{\GX_{0,1}&\CX_{0,2}\ar[l]_(0.44)s&\dots\ar[l]_(0.42)s&\CX\ar[l]_(0.42)s\ar[r]^(0.44)d&\dots\ar[r]^(0.40)d&\CX_{r-2,r}
\ar[r]^(0.46)d&\GX_{r-1,r}}\]
are smooth blowups, that is, we have a strong factorization of $\psi$.
\end{theorem}

\setcounter{section}{2}\setcounter{subsection}{2}\setcounter{theorem}{5}\setcounter{equation}{2}

\subsection{Geometric quotients}\label{wors6:sssec:quotients} It is a well known fact in Mumford's Geometric Invariant Theory that a $\C^*$-action on a polarized variety $(X,L)$  allows to define projective quotients of $X$, as follows. Let 
\[A:=\bigoplus_{m\geq 0}\HH^0(X,mL)\] 
be the {\em graded $\C$-algebra of sections of $L$}, and set, for every $\tau \in [0,\delta]\cap \Q$:
 \[A(\tau):=\bigoplus_{\substack{m\geq 0\\m\tau\in\Z}}\HH^0(X,mL)_{m\tau},\qquad \GX(\tau):=\Proj(A(\tau)).
\]
We will call the  $\GX(\tau)$'s the {\em GIT-quotients of the $\C^*$-action on $(X,L)$}.

The above construction has been extended by Bia{\l}ynicki-Birula and Swi\c{e}cicka to the construction of other (in general non projective) quotients of $X$, described in terms of semi-sections and sections of the $\C^*$-action (cf. \cite{wors6-BBS}; see also \cite{wors6-WORS3}). Their construction allows us to describe easily the open subsets $X^{\ss}(\tau)\subset X$ of semistable points on which the rational maps
$$
X\dashrightarrow \GX(\tau)
$$
are defined. 
In order to do so we need the language of {\em BB-cells}, provided by the {\em Bia{\l}ynicki-Birula decomposition} of $X$ (see \cite{wors6-BB,wors6-CARRELL}): given a fixed point component $Y\in \cY$, we denote by 
\begin{equation}\label{wors6:eq:BBcells}
X^\pm(Y):=\displaystyle\{x\in X|\,\, \lim_{t^{\pm 1}\to 0} tx\in Y\},\qquad 
B^\pm(Y):=\overline{X^\pm(Y)},
\end{equation}
the {\it BB-cells} of the action and their closures. The equalities:
\[X=\bigsqcup_{Y\in \cY}X^+(Y)=\bigsqcup_{Y\in \cY}X^-(Y)
\]
are called the {\em $+$} and {\em $-$ BB-decomposition} of $X$. 

To describe $X^{\ss}(\tau)$ for a $\C^*$-action on a polarized pair $(X,L)$, let us define
$$B^\pm_i:=\bigcup_{\mu_L(Y)=a_i}B^\pm(Y)$$
and, for notational reasons, $B^\pm_i=\emptyset$ for $i\in \Z\setminus\{0,\dots,r\}$. Then,
following \cite{wors6-BBS}, we may write:
\[
X^{\ss}(\tau)=X\setminus \big(\bigcup_{a_i<\tau}B^+_i\cup \bigcup_{a_i>\tau}B^{-}_i\big).
\]
In particular, $X^{\ss}(\tau)$ is the same for any  $\tau\in (a_i,a_{i+1})\cap \Q$, and so we set:
\begin{itemize}[itemsep=5pt,topsep=5pt]
\item $\GX_{i,i}:=\GX(\tau)$ for $\tau=a_i$; 
\item $\GX_{i,i+1}:=\GX(\tau)$ for $\tau\in(a_i,a_{i+1})\cap \Q$.
\end{itemize}
Note that $\GX_{i,i+1}$ is a geometric quotient of $X$ for every $i$, whereas $\GX_{i,i}$ is only semigeometric.  
Moreover, by construction, $$\GX_{0,0}=Y_0,\quad \GX_{r,r}=Y_r.$$ 
Furthermore, since the intersection of all the open sets $X^{\ss}(i,i+1)$ is nonempty, the varieties $\GX(i,i+1)$ are birationally equivalent. The natural birational maps among them fit in the following commutative diagram, whose diagonal arrows are contractions:
\begin{equation}
\begin{tikzcd}[
  column sep={3.4em,between origins},
  row sep={3.5em,between origins},
]
&\GX_{0,1} \arrow[rr,dashed] \arrow[rd] \arrow[dl]&&\GX_{1,2} \arrow[dl]&\dots\ \ &\GX_{r-2,r-1} \arrow[rr,dashed] \arrow[rd] &&\GX_{r-1,r}\arrow[rd] \arrow[dl]&\\
\GX_{0,0}&&\GX_{1,1}&\dots&&\dots\ \ &\GX_{r-1,r-1}&&\GX_{r,r}
\end{tikzcd}
\label{wors6:eq:GITquot}
\end{equation}
The horizontal maps are Atiyah flips; this is a known fact which goes back to the work of Thaddeus, Reid and W{\l}odarczyk (\cite{wors6-Wlodarczyk2009,wors6-Thaddeus1996,wors6-ReidFlip}).

\begin{lemma}[AM vs. FM]\label{wors6:lem:AMFM}
Let $(X,L)$ be a polarized pair, with $X$ smooth, supporting an equalized $\C^*$-action, and let $C$ be the closure of a $1$-dimensional orbit, whose sink and source are denoted by $x_-$ and $x_+$. Then $C$ is a smooth rational curve of $L$-degree  equal to $\mu_L(x_+)-\mu_L(x_-)$.
\end{lemma} 

\setcounter{section}{4}\setcounter{theorem}{0}

In this section we consider a nontrivial $\C^*$-action on a polarized pair $(X,L)$ and we  study the Chow quotient of $X$ by $\C^*$, a concept due to Kapranov who, in his seminal work \cite{wors6-Kap}, defined it and identified it with a moduli of pointed curves when $X$ is a Grassmannian endowed with an action of a complex torus. Slightly different parameter spaces for $\C^*$-invariant cycles were previously proposed  by Bia{\l}ynicki-Birula and Sommese, \cite{wors6-BBSo}, inspired by a construction by Fujiki (cf. \cite{wors6-Fj}).

Following Kapranov's definition,  we may find a nonempty open set $V\subset X$ of points $x$ such that the cycles $\overline{\C^*\cdot x}$ belong to a single homology class in $X$. By shrinking $V$, if necessary, the geometric quotient of $V$ by the action of $\C^*$ exists. Moreover, taking closures in $X$ of the orbits of points of $V$, we obtain an algebraic irreducible family of rational curves in $X$, therefore a morphism $\psi:V/\C^*\to \Chow(X)$.

\begin{definition}
With the above notation, the {\em Chow quotient} of $X$ by the action of $\C^*$ is defined as the closure:
$$
\overline{\CX}:=\overline{\psi(V/\C^*)}\subset \Chow(X).
$$
We will denote by $\CX$ the normalization of $\overline{\CX}$, and call it the {\em normalized Chow quotient} of $X$ by the $\C^*$-action. 
The pullback to $\CX$ of the universal family of $\Chow(X)$ will be denoted by $p:\cU\to \CX$, with evaluation morphism $q:\cU\to X$.
\end{definition}

The  $\C^*$-action on $X$ extends naturally to an action on $\Chow(X)$ and the corresponding universal family of cycles. This action restricts to the trivial action on $V/\C^*$, hence to the trivial action on $\CX$. Consequently we have an action on $\cU$, so that the maps $p:\cU\to\CX$, $q:\cU\to X$ are $\C^*$-equivariant, and every cycle parametrized by $\CX$ is $\C^*$-invariant. 

\setcounter{theorem}{2}

Throughout the rest of the paper we will assume that  the  $\C^*$-action is  equalized and that $X$ is smooth. Let us start by observing that in this case every $\C^*$-invariant cycle parametrized by $\CX$ is a chain of rational curves with simple normal crossings
(see \cite[Lemma 3.12]{wors6-MMW} for a similar statement for the parameter space of stable maps):

\begin{lemma}\label{wors6:lem:reduced}
Let $(X,L)$ be a polarized pair, with $X$ smooth, admitting an equalized $\C^*$-action. 
Let $z\in \CX$ be any point, and let $Z=p^{-1}(z)$ be the corresponding cycle. 
Then $Z$ is reduced, and its dual graph 
 is of type $\DA$. 
\end{lemma}
\begin{proof}
Note that $Z$ is necessarily a connected cycle meeting the sink and the source, of $L$-degree equal to the bandwidth $\delta=a_r$ of the action, by Lemma \ref{wors6:lem:AMFM}. If $Z_j$ is an irreducible component of $Z$, denoting by $z_j^-,z_j^+$ its sink and source, the equalization assumption implies that the $L$-degree of the reduced scheme of $Z_j$ is equal to $\mu_L(z_j^+)-\mu_L(z_j^-)$ (again by Lemma \ref{wors6:lem:AMFM}), and this number is of type $a_{i}-a_{j}$ for some $i,j\in\{0,\dots,r\}$. We conclude by summing on the components of $Z$.
\end{proof}

\begin{remark}\label{wors6:rem:ratcomp}
Note that the proof of Lemma \ref{wors6:lem:reduced} also implies that an irreducible component $Z_k$ of $Z=p^{-1}(z)$  is a rational curve of degree $\mu_L(z_k^+)-\mu_L(z_k^-)$, where $z_k^-,z_k^+$ denote the sink and the source of $Z_k$, respectively.  In particular, there are no components of $Z$ contained in a fixed-point component.
\end{remark}

A consequence of the reducedness given by Lemma \ref{wors6:lem:reduced} is the following statement, which can be obtained by applying standard arguments on parameter spaces:

\begin{lemma}\label{wors6:lem:flat}
Let $(X,L)$ be a polarized pair, with $X$ smooth,  with an equalized $\C^*$-action.   
Then the universal family $p:\cU\to \CX$  is flat.
\end{lemma}

\begin{proof}
Let $Z=p^{-1}(z)$ be any fiber of $p$. By Lemma \ref{wors6:lem:reduced}, the components of $Z$ are all reduced, so that $Z$ is normal at the general point of each component and, since $\CX$ is normal, we may apply \cite[Theorem~I.6.6]{wors6-kollar} to claim that $p$ is flat at the general point of any component of $Z$. We then conclude by applying \cite[Theorem~I.7.3.1]{wors6-kollar}. 
\end{proof}

\begin{remark}\label{wors6:rem:Hilb}
As a consequence of the flatness of $\cU\to \CX$, we obtain a morphism $\CX\to \Hilb(X)$ such that $\cU$ is the pullback to $\CX$ of the universal family over $\Hilb(X)$. By shrinking the open subset $V$ introduced above, we can assume that the image of $\CX\to \Hilb(X)$ is the closure of the image of the natural map from $V/\C^*$ to $\Hilb(X)$, whose normalization $\HX$ we call the {\em Hilbert quotient} of the action. On the other hand we have a natural morphism from $\Hilb(X)$ to $\Chow(X)$ (cf. \cite[Theorem I.6.3]{wors6-kollar}), so it turns out that $\CX$ coincides with the Hilbert quotient of $X$. The idea of considering a flat family of $\C^*$-invariant cycles, whose general element is the closure of a general orbit --in the complex geometric setting-- goes back to Fujiki, \cite{wors6-Fj} (see also \cite[Theorem~0.1.2]{wors6-BBSo}).
\end{remark}

\begin{remark}\label{wors6:rem:chblow}
Given an equalized $\C^*$-action on $X$, every cycle parametrized by $\CX$ has a unique lift-up to the blowup $X^\flat$  of $X$ along the sink and the source. In particular  $\CX \simeq \CX^\flat$.
\end{remark}

\setcounter{section}{6}\setcounter{theorem}{3}\setcounter{equation}{13}

\begin{proposition}\label{wors6:prop:smoothChows} Under the hypotheses of Theorem \ref{wors6:thm:main}, assume moreover that $X$ is convex. Then $\CX_{i,j}$ is smooth if $i=0$ or $j=r$. %for every $0 < j \le r$, $0 \le i <r$. 
In particular $\CX$ is smooth.
\end{proposition}

\begin{proof}
Let us prove the statement for the quotients $\CX_{0,j}$; the case of $\CX_{i,r}$ is analogous.
Note that the variety $X_{0,j}$ is by definition the pruning $X(\tau_-,\tau_+)$ for $0 < \tau_- <a_1$, $a_{j-1} < \tau_+ < j$.  

Let us consider the blowup of $X$ along the sink $Y_0$ and the source $Y_r$ of the action, $\beta: X^\flat \to X$, with exceptional divisors $Y_0^\flat$ and $Y_r^\flat$.
The indeterminacy locus of the birational modification $X_{0,j} \dashrightarrow X^\flat$ is disjoint from $Y_0^\flat$, so there exists a smooth equivariant blowdown $X_{0,j} \to X'_{0,j}$  which contracts $Y_0^\flat$ to $Y_0$. Alternatively, one may describe $X'_{0,j}$ as the pruning $X(0,\tau_+)$.
The Chow quotients $\CX_{0,j}$ and $\CX'_{0,j}$ are isomorphic (Cf. Remark \ref{wors6:rem:chblow}),
so we will show that $\CX'_{0,j}$ is smooth. To simplify notation, for the rest of the proof we will set $X':=X'_{0,j}$.
The sink of the $\C^*$-action on $X'$ is (isomorphic to) $Y_0$, and the source $Y'_r$ is a divisor  which is a birational modification of $Y_r^\flat$.

Let $C$ be the closure of a $1$-dimensional orbit of the action with sink and source $x_\pm$, and assume that $x_+\in Y'_r$. Then the weights of the $\C^*$-action on $T_{X',x_+}$ (respectively $T_{X',x_-}$) are all $0$'s and $1$'s (respectively $1$'s, $0$'s and $(-1)$'s). Then by \cite[Lemma 2.16]{wors6-WORS1} we get that the degrees of the direct summands of ${T_{X'}}_{|C}$ cannot be smaller than $-1$; thus the same holds for $\cN_{C/X'}$, and we may conclude that 
\begin{equation}\label{wors6:eq:vanish}
\HH^1(C,{T_{X'}}_{|C})=\HH^1(C,\cN_{C/X'})=0.
\end{equation}


Let us now denote by $C_g$ the closure of the general orbit of the action in $X'$, and by $[C_g]$ the corresponding point in $\Hilb(X')$; applying the above argument to $C_g$ we get that $\Hilb(X')$ is smooth at $[C_{g}]$. In particular there exists a unique irreducible component  of $\Hilb(X')$ containing it, that we denote here by $\cM$. Note that $\cM$ is a projective variety.
Let us now consider the induced $\C^*$-action on $\cM$; since $\cM$ is smooth at $[C_g]$, by \cite[Theorem 2.5]{wors6-BB} there exists a unique fixed point component  $H \subset \cM$ containing $[C_g]$. We will show that $\cM$ is smooth at every point of $H$, from which the smoothness of $H$ will follow (again by \cite[Theorem 2.5]{wors6-BB}). 

By Remark \ref{wors6:rem:Hilb}, the normalization of $H$ is $\CX'$; in particular, for every $[C]\in H$, the corresponding subscheme $C$ is a reduced tree of smooth rational curves (Lemma \ref{wors6:lem:reduced}). Furthermore, the intersection of two meeting components of $C$ is transversal, since at every fixed point $x'$ the positive and the negative part of $T_{X',x'}$ are transversal. In particular, the subscheme $C$ is a local complete intersection, hence the conormal sheaf $\cI_{C,X'}/\cI^2_{C,X'}$ is locally free, and we have a short exact sequence:
$$
\shse{\cI_{C,X'}/\cI^2_{C,X'}}{{\Omega_{X'}}_{|C}}{\Omega_C}
$$
By applying the functor $\Hom(\,\underline{\ \ }\,,\cO_C)$, we get that $\Ext^1(\cI_{C,X'}/\cI^2_{C,X'},\cO_C)$ is a quotient of $\Ext^1({\Omega_{X'}}_{|C},\cO_C)=\HH^1(C,{T_{X'}}_{|C})$. To conclude that $H$  is smooth at $[C]$ it is enough to show (see, for instance, \cite[Theorem~I.2.8]{wors6-kollar}) that \begin{equation}\label{wors6:eq:h1}
\HH^1(C,{T_{X'}}_{|C})=0.
\end{equation}
To prove (\ref{wors6:eq:h1}) we  write $C$ as $C_1 \cup C_2$, with $C_2$ irreducible and meeting the source, and set $\{p\}=C_1\cap C_2$. By (\ref{wors6:eq:vanish}), we know that $\HH^1(C,T_{X'}|_{C_2})=0$,  and so if $C=C_2$ the proof is finished. If else $C_1 \not = \emptyset$ we consider the exact sequence
\[\shse{T_{X'}|_{C_1}(-p)}{T_{X'}|_C}{T_{X'}|_{C_2}}\]
Since $X'$ is isomorphic to $X$ in a neighborhood of $C_1$, and $X$ is convex by assumption, we  have $\HH^1(C_1,T_{X'}|_{C_1}(-p))=0$ by (32) in the Proof of \cite[Lemma~10]{wors6-FP}, so  (\ref{wors6:eq:h1}) holds.
\end{proof}

\setcounter{section}{7}\setcounter{theorem}{0}

\begin{corollary}\label{wors6:cor:main}
Let $(X,L)$ be a polarized pair, with $X$ smooth, endowed with an equalized nontrivial $\C^*$-action. Then the conclusions of Theorems \ref{wors6:thm:main} and \ref{wors6:thm:convex} hold if the action satisfies the following condition:
\begin{equation}\label{wors6:eq:condition}\tag{$\star$} 
\text{$\codim(B^\pm(Y),X)>1$ for every $Y\in \cY$ of weight $a_i$ with $2 \le i \le r-2$.}
\end{equation}
\end{corollary}

\begin{thebibliography}{ORSC\etalchar{+}23b+++}

\bibitem[AO02]{AO2}
M.~Andreatta and G.~Occhetta, 
\emph{Special rays in the Mori cone of a projective variety}, 
Nagoya Math.~J.\ \textbf{168} (2002), 127--137,
\doi{10.1017/S0027763000008400}.

\bibitem[BB73]{BB}
A.~Bia{\l}ynicki-Birula, 
\emph{Some theorems on actions of algebraic groups},
Ann.\ of Math.~(2) \textbf{98} (1973), 480--497,
\doi{10.2307/1970915}.

\bibitem[BBCM02]{BBC}
A.~Bia{\l}ynicki-Birula, J.\,B.~Carrell, and W.\,M.~McGovern, 
\emph{Algebraic quotients. Torus actions and cohomology. The
adjoint representation and the adjoint action}, Encyclopaedia Math.\ Sci.,
vol.~131, 
Springer-Verlag, Berlin, 2002,
\doi{10.1007/978-3-662-05071-2}.

\bibitem[CP91]{CP}
F.~Campana and T.~Peternell, 
\emph{Projective manifolds whose tangent bundles are numerically effective}, 
Math.\ Ann.\ \textbf{289} (1991), no.~1, 169--187,
\doi{10.1007/BF01446566}.

\bibitem[Car02]{CARRELL}
J.\,B.~Carrell, 
\emph{Torus actions and cohomology}, 
in: \emph{Algebraic quotients. Torus actions and cohomology. The
  adjoint representation and the adjoint action}, pp.~83--158,
Encyclopaedia Math.\ Sci., vol.~131, Springer, Berlin, 2002,
\doi{10.1007/978-3-662-05071-2\_2}.

\bibitem[Dol03]{DolgachevGIT}
I.~Dolgachev, 
\emph{Lectures on invariant theory},
London Math.\ Soc.\ Lecture Ser., 
vol~296, 
Cambridge Univ.\ Press, Cambridge, 2003,
\doi{10.1017/CBO9780511615436}.

\bibitem[FSC24]{FS}
A.~Franceschini and L.\,E.~Sol\'{a}~Conde, 
\emph{Inversion maps and torus actions on rational homogeneous va\-ri\-e\-ties}, Geom.\ Dedicata \textbf{218} (2024), no.~1,
Paper No.~21, \href{https://doi.org/10.1007/s10711-023-00866-z}{\texttt{doi:10.1007/s10711-023-\allowbreak 00866-z}}

\bibitem[HK00]{HuKeel}
Y.~Hu and S.~Keel,
\emph{Mori dream spaces and GIT} (dedicated to William Fulton on the occasion of his 60th birthday),
Michigan Math.~J.\ \textbf{48} (2000), 331--348,
\doi{10.1307/mmj/1030132722}.

\bibitem[Ive72]{IVERSEN}
B.~Iversen, 
\emph{A fixed point formula for action of tori on algebraic varieties},
Invent.\ Math.\ \textbf{16} (1972), no.~3, 229--236,
\doi{10.1007/BF01425495}.

\bibitem[Kan17]{Kan5}
A.~Kanemitsu,
\emph{Fano $5$-folds with nef tangent bundles}, 
Math.\ Res.\ Lett.\ \textbf{24} (2017), no.~5, 1453--1475,
\doi{10.4310/MRL.2017.v24.n5.a6}.

\bibitem[Kap93]{Kap}
M.\,M.~Kapranov, 
\emph{Chow quotients of Grassmannians. {I}},
in: \emph{I.\,{M}.~{G}el'fand {S}eminar},  pp.~29--110,  Adv.\ Soviet
Math., vol.~16, Part~2,
Amer.\ Math.\ Soc., Providence, RI, 1993,
\doi{10.1090/advsov/016.2/02}.

\bibitem[Kol96]{kollar}
J.~Koll{\'a}r, 
\emph{Rational curves on algebraic varieties},
Ergeb.\ Math.\ Grenzgeb.\ (3) \textbf{32}, Springer-Verlag, Berlin, 1996, 
  \doi{10.1007/978-3-662-03276-3}. 

\bibitem[Lak82]{Lak82}
D.~Laksov, 
\emph{Notes on the evolution of complete correlations}, 
in: \emph{Enumerative geometry and classical algebraic geometry}
  ({N}ice, 1981), pp.~107--132, Progr.\ Math., vol.~24, 
Birkh\"auser Boston, Boston, MA, 1982,
\doi{10.1007/BF01585097}.

\bibitem[Lak87]{Lak87}
\bysame,
\emph{Completed quadrics and linear maps}, 
 in: \emph{Algebraic geometry, {B}owdoin, 1985} ({B}runswick, {M}aine,
  1985), pp.~371--387, Proc.\ Sympos.\ Pure Math., vol.~46, Part 2,
  Amer. Math. Soc., Providence, RI, 1987,
  \doi{10.1090/pspum/046.2/927988}.

\bibitem[Lak88]{Lak88}
\bysame, 
\emph{The geometry of complete linear maps}, 
Ark.\ Mat.\  \textbf{26} (1988), no.~2, 231--263,
\doi{10.1007/BF02386122}.

\bibitem[Mas20a]{Mas1}
A.~Massarenti, 
\emph{On the birational geometry of spaces of complete forms {I}:
  collineations and quadrics}, 
Proc.\ Lond.\ Math.\ Soc.~(3) \textbf{121} (2020), no.~6, 1579--1618,
\doi{10.1112/plms.12377}.

\bibitem[Mas20b]{Mas2}
\bysame,
\emph{On the birational geometry of spaces of complete forms {II}:
  {S}kew-forms}, 
J.~Algebra \textbf{546} (2020), 178--200,
\doi{10.1016/j.jalgebra.2019.10.047}.

\bibitem[MFK94]{MFK}
D.~Mumford, J.~Fogarty, and F.~Kirwan, 
\emph{Geometric invariant theory}, 3rd ed., Ergeb.\ Math.\ Grenzgeb.~(2), 
vol.~34,  Springer-Verlag, Berlin, 1994,
\doi{10.1007/978-3-642-57916-5}.

\bibitem[MOSW\etalchar{+}15]{MOSWW}
R.~Mu{\~n}oz, G.~Occhetta, L.\,~E.~Sol{\'a}~Conde, K.~Watanabe,
  and J.\,A.~Wi\'sniewski, 
\emph{A survey on the Campana-Peternell conjecture}, 
Rend.\ Istit.\ Mat.\ Univ.\ Trieste, \textbf{47} (2015), 127--185,
\doi{10.13137/0049-4704/11223}.

\bibitem[ORSW22]{WORS3}
 G.~Occhetta, E.\,A.~Romano, L.\,E.~Sol\'a Conde, and J.\,A.~Wi\'sniewski, 
\emph{Small modifications of Mori dream spaces arising from
  {$\mathbb{C}^*$}-actions}, 
Eur.~J.\ Math.\ \textbf{8} (2022), no.~3, 1072--1104,
\doi{10.1007/s40879-022-00540-w}.



\bibitem[ORSW23a]{WORS5}
\bysame, 
\emph{Rational homogeneous spaces as geometric realizations of birational
  transformations}, 
Rend.\ Circ.\ Mat.\ Palermo~(2) \textbf{72} (2023), no.~6, 3223--3253, 
\doi{10.1007/s12215-023-00880-w}.

\bibitem[ORSW23b]{WORS1}
\bysame, 
\emph{Small bandwidth {$\mathbb C^*$}-actions and birational geometry}, 
J.~Algebraic Geom.\ \textbf{32} (2023), no.~1, 1--57,
\doi{10.1090/jag/808}.

\bibitem[ORSW25]{WORS6}
\bysame, 
\emph{Chow quotients of $\mathbb{C}^*$-actions}, 
Int. Math. Res. Not. IMRN. \textbf{12} (2025), 
\doi{https://doi.org/10.1093/imrn/rnaf154}

\bibitem[OS26]{OS2}
G.~Occhetta and L.\,E.~Sol{\'a}~Conde, 
\emph{Characterizing rational homogeneous spaces via {$\mathbb
  C^*$}-actions}, 
Collect.\ Math.\ \textbf{77} (2026), no.~1, 161--193,
\doi{10.1007/s13348-024-00463-7}.

\bibitem[RW22]{RW}
E.\,A.~Romano and J.\,A.~Wi\'{s}niewski, 
\emph{Adjunction for varieties with a {$\mathbb{C}^*$} action}, 
Transform.\ Groups \textbf{27} (2022), no.~4, 1431--1473,
\doi{10.1007/s00031-020-09627-8}.

\bibitem[Ser06]{Ser1}
E.~Sernesi, 
\emph{Deformations of algebraic schemes},
Grundlehren math.\ Wiss.,  vol.~334, 
Springer-Verlag, Berlin, 2006.

\bibitem[Tal09]{Talpo}
M.~Talpo, 
\emph{Deformation theory}, 
Master's thesis, University of Pisa, 2009, available at \url{https://people.dm.unipi.it/talpo/stuff/tesi.pdf}.

\bibitem[Tha96]{Thaddeus1996}
M.~Thaddeus, 
\emph{Geometric invariant theory and flips}, 
J.~Amer.\ Math.\ Soc.\ \textbf{9} (1996), no.~3, 691--723,
\doi{10.1090/S0894-0347-96-00204-4}.

\bibitem[Tha99]{Thaddeus}
\bysame, 
\emph{Complete collineations revisited}, 
Math.\ Ann.\ \textbf{315} (1999), no.~3, 469--495,
\doi{10.1007/s002080050324}.

\bibitem[Vai82]{Vai82}
I.~Vainsencher, 
\emph{Schubert calculus for complete quadrics}, 
in: \emph{Enumerative geometry and classical algebraic geometry}
  ({N}ice, 1981),  pp.~199--235, Progr.\ Math., vol.~24, 
Birkh\"auser, Boston, MA, 1982,
\doi{10.1007/978-1-4684-6726-0_10}.

\bibitem[Vai84]{Vai84}
  \bysame,
  \emph{Complete collineations and blowing up determinantal ideals}, 
  Math.\ Ann.\ \textbf{267} (1984), no.~3, 417--432,
  \doi{10.1007/BF01456098}.

\end{thebibliography}

\begin{thebibliography}{10}

\bibitem{wors6-AKMW}
Dan Abramovich, Kalle Karu, Kenji Matsuki, and Jaros{\l}aw W{\l}odarczyk.
\newblock Torification and factorization of birational maps.
\newblock {\em J. Amer. Math. Soc.}, 15(3):531--572, 2002.

\bibitem{wors6-AWDuke}
Marco Andreatta and Jaros{\l}aw~A. Wi\'{s}niewski.
\newblock A note on nonvanishing and applications.
\newblock {\em Duke Math. J.}, 72(3):739--755, 1993.

\bibitem{wors6-BHK}
Hendrik B\"{a}ker, J\"{u}rgen Hausen, and Simon Keicher.
\newblock On {C}how quotients of torus actions.
\newblock {\em Michigan Math. J.}, 64(3):451--473, 2015.

\bibitem{wors6-BRUS}
Lorenzo Barban, Eleonora~A. Romano, Luis E.~Sol\'a Conde, and Stefano Urbinati.
\newblock Mori dream pairs and $\mathbb{C}^*$-actions, 2022.

\bibitem{wors6-BB}
Andrzej Bia{\l}ynicki-Birula.
\newblock Some theorems on actions of algebraic groups.
\newblock {\em Ann. of Math. (2)}, 98:480--497, 1973.

\bibitem{wors6-BBSo}
Andrzej Bia{\l}ynicki-Birula and Andrew~John Sommese.
\newblock Quotients by {${\bf C}\sp{\ast} $} and {${\rm SL}(2,{\bf C})$}
  actions.
\newblock {\em Trans. Amer. Math. Soc.}, 279(2):773--800, 1983.

\bibitem{wors6-BBS}
Andrzej Bia{\l}ynicki-Birula and Joanna \'Swi{\c e}cicka.
\newblock Complete quotients by algebraic torus actions.
\newblock In {\em Group actions and vector fields ({V}ancouver, {B}.{C}.,
  1981)}, volume 956 of {\em Lecture Notes in Math.}, pages 10--22. Springer,
  Berlin, 1982.

\bibitem{wors6-CARRELL}
James~B. Carrell.
\newblock Torus actions and cohomology.
\newblock In {\em Algebraic quotients. {T}orus actions and cohomology. {T}he
  adjoint representation and the adjoint action}, volume 131 of {\em
  Encyclopaedia Math. Sci.}, pages 83--158. Springer, Berlin, 2002.

\bibitem{wors6-CS79}
James~B. Carrell and Andrew~J. Sommese.
\newblock Some topological aspects of $\mathbb{C}^*$-actions on compact kaehler
  manifolds.
\newblock {\em Commentarii mathematici Helvetici}, 54:567--582, 1979.

\bibitem{wors6-Mas3}
Elsa Corniani and Alex Massarenti.
\newblock Complete symplectic quadrics and {Kontsevich} spaces of conics in
  {Lagrangian} {Grassmannians}.
\newblock {\em Adv. Math.}, 397:50, 2022.
\newblock Id/No 108205.

\bibitem{wors6-De}
Olivier Debarre.
\newblock {\em Higher-dimensional algebraic geometry}.
\newblock Universitext. Springer-Verlag, New York, 2001.

\bibitem{wors6-DRS}
Sandra Di~Rocco and Luca Schaffler.
\newblock Families of pointed toric varieties and degenerations.
\newblock {\em Math. Z.}, 301(4):4119--4139, 2022.

\bibitem{wors6-DolgachevHu}
Igor~V. Dolgachev and Yi~Hu.
\newblock Variation of geometric invariant theory quotients.
\newblock {\em Inst. Hautes \'{E}tudes Sci. Publ. Math.}, (87):5--56, 1998.
\newblock With an appendix by Nicolas Ressayre.

\bibitem{wors6-EH00}
David Eisenbud and Joe Harris.
\newblock {\em The geometry of schemes}, volume 197 of {\em Graduate Texts in
  Mathematics}.
\newblock Springer-Verlag, New York, 2000.

\bibitem{wors6-Fj}
Akira Fujiki.
\newblock Fixed points of the actions on compact {K}\"{a}hler manifolds.
\newblock {\em Publ. Res. Inst. Math. Sci.}, 15(3):797--826, 1979.

\bibitem{wors6-FP}
W.~Fulton and R.~Pandharipande.
\newblock Notes on stable maps and quantum cohomology.
\newblock In {\em Algebraic geometry---{S}anta {C}ruz 1995}, volume~62 of {\em
  Proc. Sympos. Pure Math.}, pages 45--96. Amer. Math. Soc., Providence, RI,
  1997.

\bibitem{wors6-Ha}
Robin Hartshorne.
\newblock {\em Algebraic geometry}.
\newblock Springer-Verlag, New York-Heidelberg, 1977.

\bibitem{wors6-Hassett}
Brendan Hassett.
\newblock Moduli spaces of weighted pointed stable curves.
\newblock {\em Adv. Math.}, 173(2):316--352, 2003.

\bibitem{wors6-IVERSEN}
Birger Iversen.
\newblock A fixed point formula for action of tori on algebraic varieties.
\newblock {\em Inv. Math.}, 16(3):229--236, 1972.

\bibitem{wors6-KSZ}
M.~M. Kapranov, B.~Sturmfels, and A.~V. Zelevinsky.
\newblock Quotients of toric varieties.
\newblock {\em Math. Ann.}, 290(4):643--655, 1991.

\bibitem{wors6-Kap}
Mikhail~M. Kapranov.
\newblock Chow quotients of {G}rassmannians. {I}.
\newblock In {\em I. {M}. {G}el'fand {S}eminar}, volume~16 of {\em Adv. Soviet
  Math.}, pages 29--110. Amer. Math. Soc., Providence, RI, 1993.

\bibitem{wors6-KKLV}
Friedrich Knop, Hanspeter Kraft, Domingo Luna, and Thierry Vust.
\newblock Local properties of algebraic group actions.
\newblock {\em Algebraische Transformationsgruppen und Invariantentheorie,
  Birkh{\"a}user Basel}, pages 63--75, 1989.

\bibitem{wors6-kollar}
J{\'a}nos Koll{\'a}r.
\newblock {\em Rational curves on algebraic varieties}, volume~32 of {\em
  Ergebnisse der Mathematik und ihrer Grenzgebiete. 3. Folge. A Series of
  Modern Surveys in Mathematics [Results in Mathematics and Related Areas. 3rd
  Series. A Series of Modern Surveys in Mathematics]}.
\newblock Springer-Verlag, Berlin, 1996.

\bibitem{wors6-LS}
A.~Losev and Y.~Manin.
\newblock New moduli spaces of pointed curves and pencils of flat connections.
\newblock {\em Michigan Math. J.}, 48:443--472, 2000.
\newblock Dedicated to William Fulton on the occasion of his 60th birthday.

\bibitem{wors6-Hue}
C{\'e}sar Lozano~Huerta.
\newblock Birational geometry of the space of complete quadrics.
\newblock {\em Int. Math. Res. Not.}, 2015(23):12563--12589, 2015.

\bibitem{wors6-Mas1}
Alex Massarenti.
\newblock On the birational geometry of spaces of complete forms. {I}:
  {Collineations} and quadrics.
\newblock {\em Proc. Lond. Math. Soc. (3)}, 121(6):1579--1618, 2020.

\bibitem{wors6-MMW}
Mateusz Micha{\l}ek, Leonid Monin, and Jaros{\l}aw~A. Wi\'sniewski.
\newblock Maximum likelihood degree and space of orbits of a
  $\mathbb{C}^*$-action.
\newblock {\em SIAM J. Appl. Algebra Geometry}, 1(5):60--85, 2021.

\bibitem{wors6-Morelli}
Robert {Morelli}.
\newblock {The birational geometry of toric varieties.}
\newblock {\em {J. Algebr. Geom.}}, 5(4):751--782, 1996.

\bibitem{wors6-MFK}
David Mumford, John Fogarty, and France Kirwan.
\newblock {\em Geometric invariant theory}, volume~34 of {\em Ergebnisse der
  Mathematik und ihrer Grenzgebiete (2) [Results in Mathematics and Related
  Areas (2)]}.
\newblock Springer-Verlag, Berlin, third edition, 1994.

\bibitem{wors6-WORS4}
Gianluca Occhetta, Eleonora~A. Romano, Luis E.~Sol\'a Conde, and Jaros{\l}aw~A.
  Wi{\'s}niewski.
\newblock Geometric realizations of birational transformations via
  $\mathbb{C}^*$-actions.
\newblock {\em Math. Z.}, 304, 2023.

\bibitem{wors6-WORS5}
Gianluca Occhetta, Eleonora~A. Romano, Luis E.~Sol\'a Conde, and Jaros{\l}aw~A.
  Wi{\'s}niewski.
\newblock Rational homogeneous spaces as geometric realizations of birational
  transformations.
\newblock {\em Rend. Circ. Mat. Palermo, II. Ser}, 72:3223--3253, 2023.

\bibitem{wors6-WORS3}
Gianluca Occhetta, Eleonora~A. Romano, Luis~E. Sol\'{a}~Conde, and Jaros\l
  aw~A. Wi\'{s}niewski.
\newblock Small modifications of {M}ori dream spaces arising from
  {$\Bbb{C}^*$}-actions.
\newblock {\em Eur. J. Math.}, 8(3):1072--1104, 2022.

\bibitem{wors6-WORS1}
Gianluca Occhetta, Eleonora~A. Romano, Luis~E. Sol\'{a}~Conde, and Jaros\l
  aw~A. Wi\'{s}niewski.
\newblock Small bandwidth {$\Bbb C^*$}-actions and birational geometry.
\newblock {\em J. Algebraic Geom.}, 32(1):1--57, 2023.

\bibitem{wors6-Pandharipande}
Raoul Pandharipande.
\newblock Convex rationally connected varieties.
\newblock {\em Proc. Amer. Math. Soc.}, 141(5):1539--1543, 2013.

\bibitem{wors6-post}
Alexander Postnikov.
\newblock Permutohedra, associahedra, and beyond.
\newblock {\em Int. Math. Res. Not. IMRN}, (6):1026--1106, 2009.

\bibitem{wors6-ReidFlip}
Miles Reid.
\newblock What is a flip? %\url
  {http://homepages.warwick.ac.uk/staff/Miles.Reid/3folds}, 1992.

\bibitem{wors6-RW}
Eleonora~A. Romano and Jaros\l aw~A. Wi\'{s}niewski.
\newblock Adjunction for varieties with a {$\Bbb{C}^*$} action.
\newblock {\em Transform. Groups}, 27(4):1431--1473, 2022.

\bibitem{wors6-Thaddeus1996}
Michael Thaddeus.
\newblock Geometric invariant theory and flips.
\newblock {\em J. Amer. Math. Soc.}, 9(3):691--723, 1996.

\bibitem{wors6-Thaddeus}
Michael Thaddeus.
\newblock Complete collineations revisited.
\newblock {\em Math. Ann.}, 315(3):469--495, 1999.

\bibitem{wors6-Wlodarczyk}
Jaros{\l}aw W{\l}odarczyk.
\newblock Birational cobordisms and factorization of birational maps.
\newblock {\em J. Algebraic Geom.}, 9(3):425--449, 2000.

\bibitem{wors6-Wlodarczyk2009}
Jaros{\l}aw W{\l}odarczyk.
\newblock Simple constructive weak factorization.
\newblock In {\em Algebraic geometry---{S}eattle 2005. {P}art 2}, volume~80 of
  {\em Proc. Sympos. Pure Math.}, pages 957--1004. Amer. Math. Soc.,
  Providence, RI, 2009.

\bibitem{wors6-Ziegler}
G\"{u}nter~M. Ziegler.
\newblock {\em Lectures on polytopes}, volume 152 of {\em Graduate Texts in
  Mathematics}.
\newblock Springer-Verlag, New York, 1995.

\end{thebibliography}
\end{document}
